\section{Experiments and executable predictive coordinates}\label{sec:experiments}
\subsection{The physical interface}
Preparation, hidden physical states, actions and reports are standard Borel, with normalized nonnegative causal kernels. Positivity means positivity preserving, not a lower bound on every report probability. A preparation is a paid physical operation. A report acquisition and a terminal audit are also paid operations. In the relative-transport experiment there are $n\ge1$ report acquisitions between one preparation and one audit, hence $N=n+2$ physical calls. The observer must commit its response before the audit. The environment-owned audit is not feedback, cannot be replaced by a simulator, and supplies no clock before that deadline.

The sequential model declares phase-dependent alphabets $S_t$, $|S_t|\le m_t$, and Borel update probabilities $k_t(j\mid s,y)$. This model is externally scheduled; it is used first to identify the local causal cost. The autonomous model instead has one finite set $S$, one reused Borel update row, an initial state, and a halt/readout map depending only on the retained state. It must halt exactly after $n$ report acquisitions with probability one. There is no free phase, history tape, real-valued posterior register or persistent random seed. Fresh randomness may be used for the current update; any part surviving to the next call belongs to the retained label. The terminal readout is a fixed function of the final retained label, with values in the declared compact response set. A transient real-valued answer wire surviving until the audit is not a free extra register. Its real-number representation, program description, temporary workspace and time are distinct resources.

\subsection{Raw finite test closure}
Let $\mathcal V$ be a finite-dimensional real space spanned by the constant test and finitely many centered, bounded, executable future test evaluations. The space is constructed by test pullbacks of the raw kernels, not by prescribing a covariance or Fisher matrix. Coordinates in $\mathcal V/\R\one$ are selected once and their Euclidean norm is fixed by the squared task. Write $p(h)\in\R^k$ for the conditional predictions after a visible history $h$. The physical states need not be finite-dimensional.

The transport hypothesis is the following identity on this test space. Reports have a fixed probability law $\lambda$, independent of the entering physical state and of all earlier reports. The positive instrument at report $y$ preserves constants and pulls back the centered tests by a Borel matrix $A(y)$, so that
\begin{equation}\label{eq:raw-pullback}
 \E[p_t\mid\mathcal H_{t-1},Y_t=y]=A(y)p_{t-1},
 \qquad p_t=A(Y_t)p_{t-1}.
\end{equation}
The second equality concerns the full visible history and follows when these tests close under the reported instrument. We require both facts as raw test identities; a state-dependent likelihood generally does not satisfy them. Stage-dependent matrices and reference laws are allowed. Throughout the general energy theorem assume
\begin{equation}\label{eq:second-isotropy}
 \int A_t(y)^TA_t(y)\,\lambda_t(dy)=\kappa_t I,
 \qquad 0<\kappa_t<\infty,
\end{equation}
and bounded square moments. The normalized zero case can be treated directly and is excluded from division by $\kappa_t$. These conditions are checked by integration on the primitive report space. They are not assumptions about an observer's proposed quantizer.

There is a performed bounded terminal vector $Z$ with $\E[Z\mid\mathcal H_n]=p_n$. Its second moment is determined by the preparation and kernels. Under squared loss, for any retained-state response $a(S_n)$,
\begin{equation}\label{eq:projection-task}
 \E\norm{Z-a(S_n)}^2
 =B_n+\E\norm{p_n-a(S_n)}^2,
 \qquad B_n=\E\norm Z^2-\E\norm{p_n}^2.
\end{equation}
The same $B_n$ is used for checkpoint and online comparisons. In the orbit experiments $\norm{p_n}=1$ and $\E\norm Z^2=V$, hence $B=V-1$. Projecting a response into the closed unit ball does not increase loss to a predictive point in that ball. The response coordinates may be unbiased estimators assembled from a randomly selected bounded physical test; they need not be jointly measured incompatible observables.

\begin{lemma}[Realized predictive quotient]\label{lem:quotient-native}
Suppose the prescribed executable tests have conditional evaluations affine in $p(h)$, contain the coordinate tests, and include their legal finite continuations. If $p$ is Borel with standard-Borel image and a Borel section, equality of $p$, together with visible interface coordinates, is precisely predictive equivalence for this test family. The quotient is minimal among Borel statistics carrying all these test predictions. Fiber-invariant report and update kernels, jointly Borel in any legal action, descend to it.
\end{lemma}
\begin{proof}
Equality of the coordinates implies equality of every affine test evaluation; the converse follows from the included coordinate tests. Any Borel statistic carrying them factors $p$ coordinatewise. Evaluation at a Borel section defines the descended kernels and fiber invariance makes their value independent of the section. Joint action measurability is retained by composition. Legal phases and available budgets are not identified merely because the numerical coordinates coincide.
\end{proof}
For the raw orbit realizations, the image is the sphere, or the sphere together with its center after a hidden transform. A Borel rotation sending a fixed unit vector to a prescribed direction can be selected on finitely many coordinate charts. A one-call visible history therefore supplies a section. The full future experiment consists of relative instruments and the declared linear audit tests, so these coordinates are its predictive quotient, not an asserted quotient of additional unmodeled physical tasks. The general result permits a task factor when other tests are explicitly left outside this family.

\subsection{The acquired orbit}
In the isotropic case $A(y)=R(y)$ is orthogonal, and the report law pushes forward to Haar probability on $SO(k)$, with $k=d+1\ge2$. The prepared predictive point $p_0=x_*$ is a known unit vector. At each call
\begin{equation}\label{eq:relative}
 p_t=R_t p_{t-1},\qquad R_1,\ldots,R_n\ \hbox{independent Haar}.
\end{equation}
Thus the actual acquired law of $p_t$ is uniform sphere measure $\sigma_d$ for every $t\ge1$. This law is derived from the performed reports. Conditional on a report the map is invertible and no preparation restarts the prediction. The unconditional transition being isotropic does not give the observer an absolute-state report.

Additional declared actions may apply known orthogonal transformations before the independent Haar report. The proof conditions on the current action as well as the old label, and Haar invariance removes that transformation. It covers Borel randomized choices made from the counted label and fresh randomness among these actions. A controller with access to additional uncharged history, or sensing actions that directly reveal $p_t$, is not included.
